% chapters/10/latihan.tex
\bagian{Latihan dan Refleksi}

\begin{enumerate}
  \item Jelaskan bagaimana sebuah peta berpetak diubah menjadi graf. Sebutkan
        apa yang menjadi simpul, apa yang menjadi sisi, dan berapa biaya setiap
        sisinya.

  \item Jelaskan perbedaan cara kerja pencarian melebar (BFS) dan algoritma
        Dijkstra. Keadaan seperti apa yang membuat BFS tidak memadai?

  \item Tuliskan rumus yang menentukan urutan pemeriksaan pada algoritma A*,
        lalu jelaskan arti setiap unsurnya.

  \item Sebuah petak berjarak enam langkah dari titik awal dan berjarak lima
        satuan Manhattan dari titik tujuan. Hitunglah nilai $f$-nya.

  \item Jelaskan syarat yang harus dipenuhi sebuah heuristik agar jalur yang
        dihasilkan tetap terpendek. Periksa apakah jarak Manhattan dan jarak
        Euclidean memenuhi syarat itu pada peta yang mengizinkan empat arah.

  \item Jelaskan mengapa jarak Manhattan tidak memenuhi syarat itu pada peta
        yang mengizinkan delapan arah, beserta akibatnya.

  \item Pada Tabel~\ref{tab:perbandingan-heuristik}, biaya jalur ketiga
        heuristik sama sedangkan banyaknya petak yang diperiksa berbeda.
        Jelaskan mengapa demikian.

  \item Jelaskan bagaimana jalur direkonstruksi dari tabel asal. Mengapa tabel
        itu perlu memuat titik awal sebagai tanda berhenti?
\end{enumerate}

\begin{refleksibox}
\begin{enumerate}
  \item Pada langkah keenam praktikum, berapa biaya jalur yang dihasilkan
        heuristik yang digandakan? Apakah Anda menyadarinya tanpa memeriksa
        tabel perbandingan?
  \item Ketika Anda membuat tujuan tidak dapat dicapai, apakah program
        menanganinya dengan baik? Menurut Anda, apa yang seharusnya dilakukan
        NPC pada keadaan itu?
  \item Peta pada bab ini berukuran 32 kali 20 petak. Menurut perkiraan Anda,
        berapa kali lipat petak yang diperiksa Dijkstra dibandingkan Manhattan
        apabila petanya sepuluh kali lebih luas? Bagaimana Anda dapat menguji
        perkiraan itu?
\end{enumerate}
\end{refleksibox}
